\section{Planificación: backlog y verificación de viabilidad}
\label{sec:planificacion}

\subsection{Naturaleza del backlog}

Conforme a la decisión metodológica de la sección~\ref{sec:metodologia}, el backlog constituye el
instrumento de planificación y medición del proyecto. Su detalle completo se presenta en el
anexo~\ref{anx:backlog}.

\begin{longtable}{@{}L{3.0cm}C{2.2cm}L{8.3cm}@{}}
\filaencab \encab{Elemento} & \encab{Cantidad} & \encab{Descripción} \\ \midrule \endhead
Épicas & 14 & Todas vinculadas a un objetivo específico \\
Historias de usuario & 25 & Con rol, capacidad y beneficio \\
Tareas técnicas & 38 & Elementos sin usuario directo asociado \\
\midrule
\textbf{Total estimado} & \textbf{63} & \textbf{275 horas} \\
Núcleo demostrativo & 33 & \textbf{165 horas} \\
\bottomrule
\caption{Composición del backlog.}
\label{tab:composicionbacklog}
\end{longtable}

\subsubsection{Criterio de estimación}

Las estimaciones se expresan en \textbf{horas ideales de trabajo efectivo} y no en puntos de historia.
La razón es la verificabilidad: el presupuesto de tiempo disponible se conoce, de modo que las horas
permiten comprobar si el alcance cabe en el plazo, comprobación que una escala abstracta no permite. La
sección~\ref{sec:viabilidad} demuestra que esa comprobación resultó determinante.

\subsubsection{Grado de refinamiento}

Los 33 elementos del núcleo demostrativo cuentan con criterios de aceptación detallados. Los restantes
se dejan deliberadamente sin refinar hasta que su ejecución se aproxime, dado que refinar por adelantado
todo el backlog constituiría trabajo perdido si el alcance se ajusta.

\subsection{Roles de usuario}

\begin{longtable}{@{}L{3.6cm}L{9.9cm}@{}}
\filaencab \encab{Rol} & \encab{Descripción} \\ \midrule \endhead
\textbf{Administrador de sitio} & Responsable técnico de la infraestructura de sensores de una organización. Registra dispositivos, vigila la salud de los enlaces y configura umbrales \\
\textbf{Usuario de consulta} & Consulta mediciones e historial sin capacidad de modificar la configuración \\
\textbf{Integrador} & Conecta un sensor o sistema externo a la plataforma; consume la interfaz de programación y la especificación del protocolo \\
\bottomrule
\caption{Roles de usuario considerados en las historias.}
\label{tab:roles}
\end{longtable}

\subsection{Épicas}

\begin{longtable}{@{}C{1.6cm}L{6.6cm}C{1.8cm}C{1.6cm}C{1.4cm}@{}}
\filaencab \encab{Épica} & \encab{Nombre} & \encab{Objetivo} & \encab{Etapa} & \encab{Horas} \\ \midrule \endhead
EP-01 & Infraestructura y entorno de desarrollo & OE5 & 2 & 14 \\
EP-02 & Diseño de arquitectura, red y datos & OE3 & 2 & 26 \\
EP-03 & Especificación formal del protocolo & OE2 & 2 & 16 \\
EP-04 & Ingesta multiprotocolo & OE5 & 2, 3 & 26 \\
EP-05 & Normalización y máquina de estados & OE5 & 3 & 19 \\
EP-06 & Persistencia y consulta & OE5 & 2, 3 & 17 \\
EP-07 & Inventario y aprovisionamiento & OE5 & 3 & 15 \\
EP-08 & Supervisión de salud de enlace & OE5 & 3 & 18 \\
EP-09 & Seguridad y multi-tenencia & OE5 & 3 & 16 \\
EP-10 & Interfaz de administración & OE5 & 3, 4 & 30 \\
EP-11 & Alertas y notificación & OE5 & 4 & 10 \\
EP-12 & Modelo de despliegue y costos & OE4 & 2 & 8 \\
EP-13 & Validación y pruebas & OE6 & 3, 4 & 26 \\
EP-14 & Documentación y entregables & Todos & 1--4 & 34 \\
\midrule
\multicolumn{4}{@{}r}{\textbf{Total}} & \textbf{275} \\
\bottomrule
\caption{Épicas del proyecto con su objetivo asociado y esfuerzo estimado.}
\label{tab:epicas}
\end{longtable}

\subsection{Definición del núcleo demostrativo}

\begin{cajaclave}[title={Criterio de demostración de la tesis}]
El prototipo demuestra la tesis si recibe eventos por \textbf{al menos dos transportes distintos}, los
normaliza a un formato único, \textbf{detecta que un dispositivo dejó de reportar}, lo presenta en una
interfaz de gestión con el estado de cada enlace, y permite \textbf{exportar los datos en formato
abierto}.
\end{cajaclave}

\begin{longtable}{@{}L{5.6cm}L{8.0cm}@{}}
\filaencab \encab{Qué demuestra} & \encab{Elementos del núcleo} \\ \midrule \endhead
Ingesta por más de un transporte & HU-01, HU-02 o HU-03, TT-16, RF-05, RF-06 \\
Normalización a formato único & TT-17, HU-05, TT-18 \\
Supervisión de salud del enlace & HU-14, HU-15, HU-16 \\
Persistencia y consulta & TT-19, HU-07 \\
Interfaz de gestión & HU-19, HU-20, HU-21 \\
Ausencia de dependencia de proveedor & HU-08 \\
Seguridad mínima & TT-20, TT-21, TT-22 \\
Extensibilidad verificable & RNF-12, validada en el diseño \\
\bottomrule
\caption{Composición del núcleo demostrativo.}
\label{tab:nucleo}
\end{longtable}

\subsection{Verificación de viabilidad del alcance}
\label{sec:viabilidad}

Las cifras de esta subsección se obtuvieron sumando programáticamente las estimaciones de los 63
elementos del backlog, y no estimando el total de forma agregada. El procedimiento es reproducible
sobre el documento fuente.

\subsubsection{Presupuesto disponible}

Cuatro semanas por etapa a aproximadamente 15 horas semanales, conforme a la restricción RE-01:
Etapa 1, 21 horas; Etapas 2, 3 y 4, 60 horas cada una. Total: \textbf{201 horas}.

\subsubsection{Contraste entre alcance y presupuesto}

\begin{longtable}{@{}C{1.6cm}C{2.4cm}C{2.4cm}C{2.2cm}C{2.4cm}C{2.2cm}@{}}
\filaencab \encab{Etapa} & \encab{Backlog completo} & \encab{Solo núcleo} & \encab{Presup.} & \encab{Holgura completo} & \encab{Holgura núcleo} \\ \midrule \endhead
1 & 10 h & 4 h & 21 h & $+$52\% & $+$81\% \\
2 & 84 h & 53 h & 60 h & \textbf{$-$40\%} & $+$12\% \\
3 & \textbf{115 h} & \textbf{97 h} & 60 h & \textbf{$-$92\%} & \textbf{$-$62\%} \\
4 & 66 h & 11 h & 60 h & $-$10\% & $+$82\% \\
\midrule
\textbf{Total} & \textbf{275 h} & \textbf{165 h} & \textbf{201 h} & \textbf{$-$37\%} & \textbf{$+$18\%} \\
\bottomrule
\caption{Contraste entre el esfuerzo estimado y el presupuesto de tiempo disponible.}
\label{tab:viabilidad}
\end{longtable}

\begin{cajahallazgo}[title={Primer hallazgo: el backlog completo no cabe en el tiempo disponible}]
El alcance completo excede el presupuesto en un \textbf{37\%} y no resulta ejecutable. El núcleo
demostrativo, en cambio, cabe con un \textbf{18\% de holgura}.

\textbf{Decisión de planificación:} el núcleo de 165 horas constituye el compromiso del proyecto; los
30 elementos restantes se declaran \textbf{alcance condicional}, ejecutables únicamente contra holgura
verificada al cierre de cada etapa.

El fundamento de esta decisión es que la pauta mide el avance como porcentaje de tareas completadas.
Comprometer 275 horas y ejecutar 165 arrojaría un cumplimiento del 60\%, por debajo del umbral de
aprobación del 70\%; comprometer 165 y ejecutar 165 arroja un 100\%. \textbf{El mismo trabajo
produce resultados opuestos según cómo se declare el alcance.}
\end{cajahallazgo}

\subsubsection{Redistribución correctiva}

\begin{cajahallazgo}[title={Segundo hallazgo: el núcleo estaba mal distribuido en el tiempo}]
Aun limitándose al núcleo, la distribución inicial resultaba defectuosa: la Etapa 3 concentraba 97 horas
frente a 60 disponibles, mientras la Etapa 4 empleaba 11 de 60. La causa fue situar la totalidad de la
implementación en la Etapa 3, cuando \textbf{la pauta exige allí el 50\% y el 100\% en la Etapa 4}.
\end{cajahallazgo}

Se trasladan de la Etapa 3 a la Etapa 4 aproximadamente 37 horas del núcleo, correspondientes a
elementos que no forman parte de la cadena mínima que debe demostrarse funcionando al cierre de la
Etapa 3: autenticación de ingesta, aislamiento entre sitios, cifrado del transporte, consulta de
historial, exportación, vista de inventario, pruebas del protocolo e indicador agregado de salud.

\begin{longtable}{@{}C{2.0cm}C{2.6cm}C{2.6cm}C{2.6cm}@{}}
\filaencab \encab{Etapa} & \encab{Núcleo} & \encab{Presupuesto} & \encab{Holgura} \\ \midrule \endhead
1 & 4 h & 21 h & $+$81\% \\
2 & 53 h & 60 h & $+$12\% \\
3 & 60 h & 60 h & 0\% \\
4 & 48 h & 60 h & $+$20\% \\
\midrule
\textbf{Total} & \textbf{165 h} & \textbf{201 h} & \textbf{$+$18\%} \\
\bottomrule
\caption{Distribución del núcleo tras la redistribución correctiva.}
\label{tab:redistribucion}
\end{longtable}

\textbf{Verificación del umbral de la Etapa 3.} La pauta exige el 50\% de las tareas de implementación
completadas. El esfuerzo total de implementación asciende a 165 horas; lo ejecutado al cierre de la
Etapa 3 corresponde a 28 horas adelantadas en la Etapa 2 más 60 horas de la Etapa 3, esto es, 88 horas
equivalentes al \textbf{53\%}. El umbral se cumple.

\subsubsection{Orden de recorte}

Los 110 horas fuera del núcleo se ordenan por valor decreciente y se ejecutan solo contra holgura
verificada.

\begin{longtable}{@{}C{1.4cm}L{6.0cm}C{1.6cm}L{4.4cm}@{}}
\filaencab \encab{Orden} & \encab{Elementos} & \encab{Horas} & \encab{Valor si se ejecuta} \\ \midrule \endhead
1 & HU-04 Receptor DC-09 nativo & 7 & Alto: es el diferenciador frente a la oferta regional \\
2 & TT-32 Validación con sensor físico & 3 & Alto: evidencia empírica de extremo a extremo \\
3 & HU-11, HU-12 Aprovisionamiento y agrupación & 8 & Medio-alto: completa la vista de administrador \\
4 & HU-24 Notificación de condición crítica & 5 & Medio \\
5 & HU-23 Vista por tecnología de red & 3 & Medio: refuerza el argumento multiprotocolo \\
6 & HU-09 Interfaz de consulta externa & 4 & Bajo para la demostración \\
7 & Resto de elementos & 80 & Variable \\
\bottomrule
\caption{Orden de ejecución del alcance condicional.}
\label{tab:recorte}
\end{longtable}

\begin{cajaaviso}[title={Límite del recorte}]
Los 33 elementos del núcleo no se recortan bajo ninguna circunstancia: su ausencia invalidaría la
demostración en lugar de reducirla. Si el tiempo no alcanzara siquiera para el núcleo, la vía correcta
no es recortar en silencio sino \textbf{renegociar el alcance con el profesor guía y documentar la
desviación}, conducta que la pauta valora explícitamente frente al incumplimiento no justificado.
\end{cajaaviso}

\subsubsection{Reserva para imprevistos}

El 18\% de holgura del núcleo, equivalente a 36 horas, constituye la única reserva. \textbf{No se
infló ninguna estimación individual}: hacerlo habría producido un plan con apariencia de viabilidad sin
mejorarla, y habría impedido detectar la sobrecarga que esta subsección documenta. El riesgo R4 se
actualiza en consecuencia, de probabilidad media a confirmado.

\subsection{Criterio de priorización}

El orden de ejecución no sigue la numeración de las épicas, sino tres reglas:

\begin{enumerate}
\item \textbf{Primero lo que desbloquea:} el entorno de desarrollo y el esquema de datos condicionan
todo lo demás.
\item \textbf{Después lo que demuestra la tesis:} la cadena ingesta, normalización, supervisión e
interfaz. Es el recorrido que un evaluador seguirá durante una demostración.
\item \textbf{Al final lo que mejora el producto:} notificación avanzada, receptor DC-09, interfaz de
consulta externa. Valiosos, pero su ausencia no invalida la demostración.
\end{enumerate}

\begin{cajaclave}[title={Regla de decisión ante retraso}]
Se sacrifica amplitud antes que profundidad. Resulta preferible una cadena completa con dos transportes
y supervisión operando de extremo a extremo, que cuatro transportes a medio terminar: \textbf{un
recorrido completo demuestra la hipótesis; cuatro recorridos parciales no demuestran ninguna.}
\end{cajaclave}
